%% 自動生成（checks/物理第5巻/巻末/mknext_p5.py）。各章の「工学の駅」（\phstation）の科目を集めたもの。第5巻はシリーズの最後の巻なので、後ろの巻の章のかわりに、この先の工学の科目を並べる
\begin{longtblr}{colspec={Q[l,m,wd=62mm] X[l,m]}, width=\linewidth, rowhead=1, row{1}={bg=bkVol!14, font=\sffamily\bfseries\footnotesize},
  hlines={0.4pt,bkLine}, rowsep=1.5pt, cells={font=\footnotesize}}
この巻の章 & この道具を使う工学の科目 \\
{\bknextpill{原-1}\ 原子の中をのぞく：電子と原子核の発見} & 電子デバイス・真空工学・分析化学・加速器工学 \\
{\bknextpill{原-2}\ 光はつぶでもある：光電効果と光子} & 光エレクトロニクス・計測工学・放射測温・量子光学 \\
{\bknextpill{原-3}\ 電子は波でもある：ド・ブロイ波と不確定性} & 量子力学・電子顕微鏡学・材料評価・半導体工学 \\
{\bknextpill{原-4}\ 原子が出す光：エネルギー準位とスペクトル} & 原子物理学・分光分析・X線分析・レーザー工学 \\
{\bknextpill{原-5}\ 波で書く電子：シュレディンガー方程式} & 量子力学・半導体デバイス・光デバイス・量子化学 \\
{\bknextpill{原-6}\ 壁をすりぬける：トンネル効果と調和振動子} & 半導体デバイス・表面科学・分子分光学・超伝導エレクトロニクス \\
{\bknextpill{原-7}\ 電子の席の決まり：原子と周期表} & 材料工学・量子化学・磁性材料・医用画像 \\
{\bknextpill{原-8}\ 固体の中の電子：エネルギーバンド} & 固体物理・材料工学・電子材料・超伝導工学 \\
{\bknextpill{原-9}\ 一方通行の電気：pn接合とLED・太陽電池} & 電子デバイス・太陽光発電・光通信・照明工学 \\
{\bknextpill{原-10}\ スイッチと増幅：トランジスタと電子回路} & 電子回路・集積回路・パワーエレクトロニクス・計測工学 \\
{\bknextpill{原-11}\ 質量はエネルギー：相対論と原子核} & 原子核工学・加速器工学・医学物理・宇宙工学 \\
{\bknextpill{原-12}\ こわれる原子核：放射性崩壊と半減期} & 放射線計測・核医学・放射化学・年代学 \\
{\bknextpill{原-13}\ 放射線が物にあたる：放射線と物質} & 放射線工学・医学物理・遮へい設計・非破壊検査 \\
{\bknextpill{原-14}\ 放射線をはかる・身を守る：計測と防護} & 放射線管理・保健物理・医学物理・非破壊検査 \\
{\bknextpill{原-15}\ 原子炉の中の中性子：核分裂と原子炉} & 原子炉物理・原子炉工学・放射線工学・核融合工学 \\
{\bknextpill{原-16}\ 原子の力とつきあう} & 原子力工学・放射線工学・保健物理学・エネルギー工学 \\
\end{longtblr}
